\section{The Language Problem: Where Deep Learning Struggled}

\subsection{Language is Different}

Images have a neat, grid-like structure.
Language does not.

A sentence is a sequence --- each word depends on what came before it, sometimes even on what came dozens of words earlier.
The meaning of ``bank'' in ``river bank'' is completely different from ``bank'' in ``I went to the bank.''
Understanding language requires holding context across time, resolving ambiguity, and tracking long-range dependencies.

Deep learning researchers knew this was a harder problem.
They needed a new kind of architecture.

\subsection{Recurrent Neural Networks}

The first serious attempt was the \textbf{Recurrent Neural Network} (RNN).
The key difference from a standard network: RNNs have a \emph{hidden state} --- a memory that carries information forward as the network processes a sequence one element at a time.

\begin{tcolorbox}[colback=green!5!white, colframe=green!40!black, title=How an RNN Reads a Sentence]
An RNN reads words one at a time, left to right.
At each step, it updates its hidden state --- a fixed-size vector that is meant to summarize everything it has seen so far.
When it reaches the end of the sentence, the hidden state is supposed to capture the full meaning.
\end{tcolorbox}

RNNs showed promise on short sequences.
But they had a fundamental weakness: \textbf{the vanishing gradient problem}.

When backpropagation runs through a long sequence, the error signal has to travel backwards through every time step.
At each step, it gets multiplied by numbers that are often less than one.
After fifty steps, the signal has essentially vanished.
The network forgets what it read at the beginning of a long sentence by the time it reaches the end.

\subsection{LSTMs: A Smarter Memory}

In 1997, Sepp Hochreiter and Jürgen Schmidhuber introduced the \textbf{Long Short-Term Memory} (LSTM) network \cite{hochreiter1997long}.
LSTMs added a more sophisticated memory mechanism --- explicit \emph{gates} that control what information to keep, what to forget, and what to output at each step.

\begin{tcolorbox}[colback=blue!5!white, colframe=blue!40!black, title=LSTM Gates in Plain English]
Think of an LSTM's memory like a whiteboard with a careful editor. \\
\textbf{Forget gate:} Decides what old information to erase. \\
\textbf{Input gate:} Decides what new information to write. \\
\textbf{Output gate:} Decides what to read off the board for the current step.
\end{tcolorbox}

LSTMs were a significant improvement.
They powered many early breakthroughs in machine translation, speech recognition, and text generation throughout the 2010s.
Google Translate switched to an LSTM-based model in 2016 \cite{wu2016google} and saw quality improvements that surprised even its own engineers.

\subsection{Still, Fundamental Limits}

But even LSTMs had limits.
They processed sequences one word at a time --- which meant they could not be parallelized efficiently.
Training on large datasets was slow.
And the fixed-size hidden state remained a bottleneck: no matter how smart the gating mechanism, compressing an entire paragraph into a single vector loses information.

Researchers experimented with different architectures: bidirectional RNNs (reading forward and backward), encoder-decoder models, attention mechanisms bolted onto RNNs \cite{bahdanau2014neural}.
Progress was steady but incremental.

The field was working hard at a difficult problem, and it was clear that a more fundamental rethinking was needed.
